% sections/07-uml-diagrams.tex — UML Diagrams and Design Documentation

\section{UML Diagrams and Design Documentation}

% ── Requirement ──
\subsection{Requirement}

You must create UML diagrams to document your software design. These diagrams should accurately represent the structure and interactions of your system.

% ── Required Diagrams ──
\subsection{Required Diagrams}

\subsubsection{Class Diagram}

A class diagram shows the static structure of your system---its classes, attributes, methods, and relationships. Your class diagram must include:

\begin{itemize}
  \item All major classes in your project
  \item Public attributes and methods
  \item \textbf{Inheritance} --- solid line with hollow triangle arrowhead
  \item \textbf{Composition} --- solid line with filled diamond (strong ownership)
  \item \textbf{Aggregation} --- solid line with hollow diamond (weak ownership)
  \item \textbf{Associations} --- plain solid lines
  \item \textbf{Multiplicity} --- \code{1}, \code{0..1}, \code{*}, \code{1..*}
\end{itemize}

\subsubsection{Sequence Diagram (Optional but Recommended)}

A sequence diagram shows how objects interact over time. It is especially useful for documenting:

\begin{itemize}
  \item Main application flow
  \item Complex operations involving multiple objects
  \item Request/response patterns
\end{itemize}

\subsubsection{Component Diagram (Optional)}

A component diagram shows the high-level organization of your modules and their dependencies. Useful for understanding the overall architecture.

% ── UML Tools ──
\subsection{UML Tools}

\textbf{Free / Open Source:}
\begin{itemize}
  \item \href{https://plantuml.com/}{PlantUML} --- text-based, version-control friendly
  \item \href{https://app.diagrams.net/}{Draw.io} --- web-based, intuitive drag-and-drop
  \item \href{https://umbrello.kde.org/}{Umbrello} --- KDE UML modeller
  \item \href{https://staruml.io/}{StarUML} --- multi-platform UML tool
  \item \href{https://mermaid.js.org/}{Mermaid} --- text-based, renders in Markdown
\end{itemize}

\textbf{Commercial:}
\begin{itemize}
  \item Visual Paradigm
  \item Enterprise Architect
  \item Lucidchart
\end{itemize}

% ── PlantUML Example ──
\subsection{PlantUML Example}

\begin{lstlisting}[style=plaintext]
@startuml

class LinkedList<T> {
    - head : Node<T>*
    - tail : Node<T>*
    - size : int
    + insert(index : int, value : T) : void
    + remove(index : int) : T
    + get(index : int) : T
    + isEmpty() : bool
    + getSize() : int
}

class Node<T> {
    + data : T
    + next : Node<T>*
}

LinkedList *-- Node : contains

@enduml
\end{lstlisting}

% ── Diagram Placement ──
\subsection{Diagram Placement}

Store all diagrams and design documents in the \filepath{docs/} directory:

\begin{lstlisting}[style=tree]
docs/
|-- uml/
|   |-- class_diagram.png
|   |-- class_diagram.puml ......... PlantUML source
|   +-- sequence_main_flow.png
|-- design_document.md
+-- Doxyfile
\end{lstlisting}

% ── Design Document ──
\subsection{Design Document}

Create a \filename{design\_document.md} in your \filepath{docs/} directory. It should include:

\begin{enumerate}
  \item \textbf{Architecture Overview} --- high-level description of the system
  \item \textbf{Module Descriptions} --- purpose and responsibility of each module
  \item \textbf{Data Structure Choices} --- which data structures you chose and why
  \item \textbf{Algorithm Explanations} --- key algorithms with time/space complexity
  \item \textbf{Design Decisions} --- trade-offs you considered and your rationale
\end{enumerate}

% ── UML Notation Quick Reference ──
\subsection{UML Notation Quick Reference}

\begin{center}
\begin{tabular}{l l}
\hline
\rowcolor{tableheader} \textbf{Relationship} & \textbf{Notation} \\
\hline
Inheritance & Solid line with hollow triangle \\
Interface Implementation & Dashed line with hollow triangle \\
Composition & Solid line with filled diamond \\
Aggregation & Solid line with hollow diamond \\
Association & Plain solid line \\
Dependency & Dashed line with arrow \\
\hline
\end{tabular}
\end{center}
